% !TeX root = ../main.tex

\chapter{Agentic AI e Design MCP (Model Context Protocol)}
\label{cap:agentic_operations_mcp}

Il cuore AI del progetto[cite: 58, 59]. Qui viene spiegata la netta separazione tra automazione deterministica e intelligenza agentica[cite: 59, 60].

\section{Il Limite dell'Automazione Deterministica}
\label{sec:limite_automazione}

Normalmente l'assegnazione degli ordini ai vettori verrebbe gestita tramite un approccio deterministico basato su alberi decisionali (scripting \textit{if-else}). Sebbene questo metodo sia efficiente per carichi di lavoro lineari, si rivela inadeguato nel dominio complesso affrontato dal nostro sistema. 

Nel nostro digital twin, ogni missione è influenzata da una combinazione di variabili dievrse: la priorità dell'ordine, il peso del carico, la percentuale di batteria residua, la distanza dall'obiettivo e il tasso di usura del mezzo. Scrivere regole rigide per coprire ogni caso possibile porterebbe a un codice fragile e difficilmente manutenibile. 
L'Intelligenza Artificiale Agentica supera questo limite sostituendo la rigidità deterministica con un ragionamento probabilistico. L'Agente valuta dinamicamente i trade-off operativi (es. rischio di schianto vs urgenza) e ottimizza il matching Ordine-Drone in tempo reale, adattandosi a contesti imprevisti senza necessità di pre-programmazione esplicita.

\section{Architettura MCP}
\label{sec:architettura_mcp}
Per garantire che l'IA operi all'interno di confini sicuri, il sistema implementa un'architettura basata sul \textit{Model Context Protocol} (MCP). L'MCP funge da strato di astrazione tra i loop di ragionamento degli LLM (confinati in \texttt{logistic\_ai\_brain.py}) e l'infrastruttura fisica simulata su Kubernetes. 
L'LLM non possiede credenziali dirette o permessi di root sui database; comunica esclusivamente tramite richieste HTTP standardizzate al server MCP, il quale espone due categorie di \textit{Tools} rigorosamente tipizzati:

\begin{itemize}
    \item \textbf{Observer Tools (Read-Only):} Strumenti sicuri utilizzati dall'agente per la fase di raccolta delle informazioni necessarie. Includono \texttt{get\_drones\_status} (interrogazione API Kubernetes), \texttt{get\_drones\_telemetry} e \texttt{get\_pending\_orders} (estrazione dati in time-window da InfluxDB).
    \item \textbf{Remediation Tools (Write):} Strumenti attivi che modificano lo stato del sistema. Includono \texttt{send\_mqtt\_command} per l'assegnazione delle missioni e \texttt{scale\_drone\_deployment} per l'Auto-Scaling dei pod. Essendo potenzialmente pericolosi, la loro esecuzione all'interno del modulo \texttt{drone\_mcp\_layer.py} è subordinata a policy che vedremo successivamente
\end{itemize}

\section{Ruolo 1: The Health Agent}
\label{sec:health_agent}
L'\texttt{HealthAgent} si occupa della resilienza dell'infrastruttura, della diagnostica della flotta e del dimensionamento elastico delle risorse fisiche. L'agente opera analizzando i flussi di telemetria e agisce come un supervisore automatizzato della salute del sistema, gestendo monitoraggio, triage dell'usura e auto-scaling.

\subsection*{Monitoraggio dell'Usura e Manutenzione Autonoma}
A differenza di un sistema statico che attende il guasto hardware per intervenire, l'\texttt{HealthAgent} applica un modello di manutenzione preventiva basato sul campionamento del degrado meccanico (\textit{wear}) restituito dall'MCP. 
Quando l'agente rileva che un vettore ha raggiunto o superato la soglia critica del 95\%, isola il mezzo, traducendo il suo stato in \texttt{MAINTENANCE}.

\subsection*{Rilevamento Anomalie}
Il ciclo di ragionamento dell'agente include l'analisi logica per l'identificazione di guasti (es. esaurimento improvviso della batteria in volo o schianti). Se un drone si trova in stato di \texttt{MAINTENANCE} ma le sue coordinate spaziali differiscono da quelle dell'HUB centrale \texttt{[0.0, 0.0]}, l'\texttt{HealthAgent} deduce autonomamente che il mezzo è precipitato sul campo e innesca le procedure di recupero. 

Inoltre, l'agente analizza in modo proattivo il rapporto tra gli ordini pendenti e la capacità attuale della flotta. Se rileva un collo di bottiglia, decide a quante unità, e di conseguenza pod Kubernetes, scalare il cluster per mantenere inalterati gli SLO. In alcuni casi che vedremo successivamente richiede direttamente l'intervento umano.

\section{Ruolo 2: The Logistic Agent}
\label{sec:logistic_agent}
Il \texttt{LogisticAgent} si occupa della logica di business del sistema, gestendo la coda degli ordini asincroni inviati dal client simulator e coordinando l'effettiva assegnazione delle missioni sui singoli mezzi. Questo ruolo beneficia fortemente del paradigma agentico LLM rispetto all'automazione classica, in quanto richiede un'ottimizzazione complessa basata su vincoli fisici non lineari.

\subsection*{Logica di Assegnazione tramite Ragionamento Multi-Criterio}
L'agente logistico incrocia i dati degli ordini con lo stato telemetrico real-time fornito dai droni che si trovano in stato \texttt{IDLE}. L'assegnazione dell'ordine non avviene secondo criteri sequenziali standard (es. FIFO), ma attraverso una ponderazione del rischio:
\begin{enumerate}
    \item \textbf{Analisi dell'ordine:} Il peso del pacco (\texttt{weight\_kg}) influisce direttamente sulle performance del drone: un pacco pesante riduce la velocità di crociera e incrementa il consumo di batteria.
    \item \textbf{Contesto Spaziale:} La distanza euclidea impone un calcolo preventivo dell'autonomia necessario al viaggio.
    \item \textbf{Stato del Mezzo:} Il livello di batteria attuale e il grado di usura meccanica (\textit{wear}) accumulato.
\end{enumerate}

Se un ordine presenta priorità elevata (\textit{high priority}) ma un peso prossimo al limite massimo di 5.0 kg, l'agente eviterà tassativamente di assegnarlo a droni non in condizioni ottimali (es. $>80\%$) o con poca batteria, riducendo le probabilità di malfunzionamenti. Una volta cercata la coppia ottimale Ordine-Drone, l'agente lancia il comando MCP \texttt{send\_mqtt\_command} scatenando la partenza.

\section{Orchestrazione Dinamica}
\label{sec:logistic_ai_brain}

Il coordinamento e l'attivazione degli agenti sono gestiti dall'orchestratore centrale (\texttt{logistic\_ai\_brain.py}). Per garantire la massima efficienza economica, (\textit{Budget Compliance}) di cui parleremo in seguito, l'orchestratore utilizza un'architettura reattiva, basata sul contesto.

Si articola in tre meccanismi chiave implementati in \texttt{logistic\_ai\_brain.py}:

\begin{enumerate}
    \item \textbf{State Hashing:} Per evitare che gli agenti sprechino token LLM analizzando un sistema che non ha subito variazioni, l'orchestratore esegue richieste HTTP non agentiche per recuperare una fotografia del sistema (\texttt{fetch\_global\_state}) e la confronta con il ciclo precedente. Se l'hash coincide, il sistema bypassa completamente l'attivazione dell'IA e va in standby, risparmiando decine di migliaia di token orari.
    
    \item \textbf{Triage Manager e Context Injection:} Gli agenti non sprecano più interazioni inutili. È l'orchestratore che agisce da router (\texttt{triage\_manager}): analizza la fotografia del sistema e decide se serve l'intervento dell'Health Agent, del Logistic Agent, o di entrambi. Una volta stabilito il target, l'orchestratore filtra i dati (fornendo alla logistica, ad esempio, \textit{solo} i droni \texttt{IDLE}) e li inietta direttamente nel prompt iniziale dell'agente (\textit{Context Injection}). Questo costringe l'LLM a passare immediatamente alla fase di invocazione dei tool di rimediazione.
    
    \item \textbf{Esecuzione Parallela Asincrona:} Qualora il \textit{Triage Manager} stabilisca che sia necessaria sia un'azione di manutenzione/scaling che un'assegnazione di ordini, i due agenti non vengono più eseguiti in modo sequenziale. L'orchestratore li istanzia e li avvia parallelamente su processi separati (\texttt{threading.Thread}), facendoli successivamente ricongiungere. Questo approccio asincrono dimezza la latenza del ciclo decisionale, permettendo al sistema di rispondere in modo molto più reattivo a picchi di traffico imprevisti.
\end{enumerate}